\section{3D viewing transformation}
\label{sec:view}
\here{3}
\concept{The view matrix moves the world into the camera's frame. It is built from the eye $\vect e$, the point
looked at $\vect a$ and an up hint $\vect{up}$, by an orthonormal basis and a translation.}
\begin{equation}
\vect n=\frac{\vect e-\vect a}{\lVert\vect e-\vect a\rVert},\quad
\vect u=\frac{\vect{up}\times\vect n}{\lVert\vect{up}\times\vect n\rVert},\quad
\vect v=\vect n\times\vect u,\qquad
V=\begin{bmatrix}u_x&u_y&u_z&-\vect u\cdot\vect e\\v_x&v_y&v_z&-\vect v\cdot\vect e\\n_x&n_y&n_z&-\vect n\cdot\vect e\\0&0&0&1\end{bmatrix}
\label{eq:view}
\end{equation}
The code writes the same matrix with $\vect f=-\vect n$ (forward), $\vect s=\vect u$ (side) and $\vect u=\vect v$
(\code{math/Math3D.cpp lookAt}); the camera looks down its $-z$ axis.

\begin{figure}[H]
\centering
\begin{subfigure}{0.325\linewidth}\raw{play/play_first_person}{\linewidth}\caption{Soldier, first person.}\end{subfigure}\hfill
\begin{subfigure}{0.325\linewidth}\raw{play/play_third_person}{\linewidth}\caption{Third person (spring arm).}\end{subfigure}\hfill
\begin{subfigure}{0.325\linewidth}\raw{v5/uav_camera_view}{\linewidth}\caption{UAV-mounted gimbal camera.}\end{subfigure}
\caption{Three cameras on the same scene; each is just a different $(\vect e,\vect a,\vect{up})$. Notice in (c) that
the UAV camera's eye and axes come from the gimbal nodes of Figure~\ref{fig:sg}: the view matrix is the inverse of
the camera node's world matrix.}
\label{fig:cams}
\end{figure}

\begin{figure}[H]
\centering\wb{View_Projection}{0.95\linewidth}
\caption{Workbook sheet \code{View\_Projection}: $\vect u,\vect v,\vect n$ from cross products for the default orbit
camera, the view matrix, the perspective matrix, and the depth table used in Section~\ref{sec:projective}.}
\label{fig:wb-view}
\end{figure}

\textbf{Camera system.} Orbit (default), free fly (\key{Home}), chase, UAV camera, walk, first and third person on
any soldier (\key{Tab}, \key{V}), UAV chase / FPV / gimbal, any moving object's point of view (\key{F4}) and the
cinematic director (\key{K}). Changes of camera blend over 0.6\,s on a critically damped spring (\code{CameraRig::kBlendTime}).
\where{src/math/Math3D.cpp lookAt, src/scene/Camera.cpp compute(), src/scene/CameraRig.cpp, src/scene/UavView.cpp}
\critique{When $\vect{up}\parallel\vect n$ (looking straight down) the cross product vanishes. The top view of
Section~\ref{sec:proj} avoids this by using $\vect{up}=-\vect z$; the free camera clamps its pitch to $\pm85^\circ$.}
